% !TeX program = xelatex
% Compile: latexmk -xelatex main-cn.tex
% Compatibility shim: paper.cls contains pdfTeX-only output/font-map commands.
% XeLaTeX uses fontspec in preamble-cn.tex instead; leave paper.cls unchanged.
\RequirePackage{iftex}
\ifXeTeX
  \newcount\pdfoutput
  \def\pdfmapline#1{}
  \PassOptionsToPackage{protrusion=false,expansion=false}{microtype}
  \AddToHook{package/microtype/after}{%
    \renewcommand{\DisableLigatures}[2][]{}%
  }
\else
  \errmessage{main-cn.tex requires XeLaTeX. Run latexmk -xelatex main-cn.tex}
\fi
\documentclass[twocolumn]{paper}
% Load after paper.cls. Chinese text uses TeX Live's bundled Fandol fonts.
\usepackage[UTF8,fontset=fandol,scheme=plain,heading=false]{ctex}
% Jilin University theme overrides; paper.cls only renames the theme color.
% geometry also uses addtolist; preserve the class metadata helper.
\let\jlusavedaddtolist\addtolist
\geometry{a4paper,headheight=43pt}
\let\addtolist\jlusavedaddtolist
\DeclareTextFontCommand{\textbf}{\bfseries}
\paperlogo{figures/jlu.pdf}
\fancypagestyle{firststyle}{
  \fancyhf{}
  \fancyfoot[C]{\thepage}
  \fancyhead[L]{{\frontmedium\small\color{jlublue}JILIN UNIVERSITY}}
  \fancyhead[R]{%
    \ifdefempty{\paperlogofile}{}{\includegraphics[height=13.5mm]{\paperlogofile}}%
  }
}

% Add paper-specific packages and commands here.
\usepackage{amsmath,amssymb,amsthm}
\usepackage{enumitem}
\usepackage{tikz}
\graphicspath{{figures/}}
\setlist[itemize]{leftmargin=*,nosep}
\newtheorem{theorem}{Theorem}[section]
\newtheorem{lemma}[theorem]{Lemma}
\theoremstyle{definition}
\newtheorem{definition}[theorem]{Definition}
\theoremstyle{remark}
\newtheorem{remark}[theorem]{Remark}
\newcommand{\vect}[1]{\bm{#1}}
\newcommand{\R}{\mathbb{R}}

% Highlight box, matching the blue palette of the title and abstract.
\newtcolorbox{highlightbox}[1][Key idea]{
  colback=jlublue!4,colframe=jlublue,boxrule=0.5pt,
  arc=1.5mm,title=#1,fonttitle=\bfseries,
  left=2mm,right=2mm,top=2mm,bottom=2mm
}

\usetikzlibrary{positioning}


% Replace the pdfTeX-specific Latin font setup locally for XeLaTeX.
\setmainfont{Latin Modern Roman}
\setsansfont{ByteSans-Medium.ttf}[Path=fonts/,BoldFont=ByteSans-Bold.ttf]
\newfontfamily\jlufrontbold{ByteSans-Bold.ttf}[Path=fonts/,BoldFont=ByteSans-Bold.ttf]
\newfontfamily\jlufrontmedium{ByteSans-Medium.ttf}[Path=fonts/,BoldFont=ByteSans-Bold.ttf]
\renewcommand{\frontsans}{\jlufrontbold\heiti}
\renewcommand{\frontmedium}{\jlufrontmedium\heiti}
\setlength{\parindent}{2em}

\renewcommand{\beginabstract}{\abstractfont\frontsans\bfseries\jlublue{摘要}\par}
\addto\captionsenglish{%
  \renewcommand{\refname}{参考文献}%
  \renewcommand{\figurename}{图}%
  \renewcommand{\tablename}{表}%
  \renewcommand{\proofname}{证明}%
}
\patchcmd{\theorem}{Theorem}{定理}{}{}
\patchcmd{\lemma}{Lemma}{引理}{}{}
\patchcmd{\definition}{Definition}{定义}{}{}
\patchcmd{\remark}{Remark}{注}{}{}
\crefname{equation}{式}{式}
\Crefname{equation}{式}{式}
\crefname{figure}{图}{图}
\Crefname{figure}{图}{图}
\crefname{table}{表}{表}
\Crefname{table}{表}{表}
\crefname{section}{节}{节}
\Crefname{section}{节}{节}
\crefname{theorem}{定理}{定理}
\Crefname{theorem}{定理}{定理}
\fancypagestyle{firststyle}{
  \fancyhf{}
  \fancyfoot[C]{\thepage}
  \fancyhead[L]{{\frontmedium\small\color{jlublue}吉林大学\quad JILIN UNIVERSITY}}
  \fancyhead[R]{%
    \ifdefempty{\paperlogofile}{}{\includegraphics[height=13.5mm]{\paperlogofile}}%
  }
}


% 替换为真实的题目、作者、学院及联系方式。
\title{中文论文标题：\\研究方法与应用示例}
\author[1]{第一作者}
\author[1]{第二作者}
\author[1\corrauthor]{第三作者}
\affiliation[1]{吉林大学某某学院，吉林长春，130012}
% 多单位示例：\affiliation[2]{合作单位，城市，邮编}
% 对应作者：\author[1,2]{作者姓名}
\contribution[\mathsection]{通讯作者}
\checkdata[通讯邮箱]{\email{author@example.com}}
\checkdata[项目地址]{\url{https://example.com/project}}
\checkdata[关键词]{吉林大学；科研论文；研究方法；实验评估}
\paperlogo{figures/jlu.pdf}

\abstract{%
此处填写中文摘要。首先介绍研究问题及其实际意义，然后概述现有方法的不足，
说明本文提出的方法及其主要特点，最后总结实验结果和研究结论。
摘要应独立、完整，避免使用未经解释的缩写，并以实际研究内容替换本段示例。
}

\begin{document}
\maketitle

\section{引言}
介绍研究背景、问题的重要性以及应用场景，说明现有方法与研究目标之间的差距。
使用参考文献支持具体论述，例如引用排版工具的经典资料~\citep{lamport1994latex}。

\paragraph{主要贡献。}
本文的研究贡献可以概括为：
\begin{itemize}
  \item 明确研究问题及其关键挑战。
  \item 提出解决问题的方法，并说明核心设计。
  \item 通过对比实验和消融实验验证方法的有效性。
\end{itemize}

\section{相关工作}
按照研究主题梳理已有工作，讨论不同方法的特点和适用条件，并说明本文方法与已有工作的关系。
可以同时引用多篇文献~\citep{lamport1994latex,knuth1984texbook}。

\section{研究方法}
\subsection{问题定义}
设输入为 $\vect{x}_i\in\R^d$，对应目标为 $y_i$。
给定数据集 $\mathcal{D}=\{(\vect{x}_i,y_i)\}_{i=1}^{N}$，优化目标为
\begin{equation}
  \mathcal{L}(\theta)=\frac{1}{N}\sum_{i=1}^{N}
  \ell\bigl(f_\theta(\vect{x}_i),y_i\bigr).
  \label{eq:objective}
\end{equation}
其中，$\theta$ 表示模型参数，$f_\theta$ 表示模型，$\ell$ 表示损失函数。

\subsection{方法设计}
详细描述方法的组成、执行步骤和实现条件。
优化目标见\cref{eq:objective}，整体流程见\cref{fig:pipeline}。
\begin{highlightbox}[核心思想]
此处填写方法的关键思想、研究假设或主要设计原则。
\end{highlightbox}

\begin{figure}[t]
  \centering
  \begin{tikzpicture}[node distance=4mm,
    stage/.style={draw=jlublue,fill=jlublue!5,rounded corners,
      minimum width=20mm,minimum height=10mm,font=\small}]
    \node[stage] (input) {输入数据};
    \node[stage,right=of input] (method) {研究方法};
    \node[stage,right=of method] (output) {输出结果};
    \draw[->,thick,color=jlublue] (input) -- (method);
    \draw[->,thick,color=jlublue] (method) -- (output);
  \end{tikzpicture}
  % 外部图片示例：\includegraphics[width=\linewidth]{your-figure.pdf}
  \caption{方法流程示例，请替换为实际研究流程。}
  \label{fig:pipeline}
\end{figure}

\begin{theorem}[示例命题]
若损失函数 $\ell$ 非负，则 $\mathcal{L}(\theta)\geq 0$。
\end{theorem}
\begin{proof}
\cref{eq:objective}中的各项均为非负数，因此其平均值非负。
\end{proof}

\section{实验与分析}
\subsection{实验设置}
说明数据集、对比方法、评价指标、参数配置、软件版本与硬件环境，
提供足够的信息以便其他研究者复现实验。

\subsection{实验结果}
\Cref{tab:results}展示三线表的写法。表中数据仅为排版示例，不是实际实验结果。
\begin{table}[t]
  \centering
  \caption{实验结果示例（数据为占位值）。}
  \label{tab:results}
  \begin{tabular}{lrr}
    \toprule
    方法 & 得分 $\uparrow$ & 耗时 $\downarrow$ \\
    \midrule
    基线方法 A & 72.1 & 1.20 \\
    基线方法 B & 75.3 & 1.05 \\
    本文方法 & \textbf{79.4} & \textbf{0.92} \\
    \bottomrule
  \end{tabular}
\end{table}
分析实验结果、误差来源和失效情形。跨双栏图表可使用 \texttt{figure*} 或 \texttt{table*}。

\subsection{消融实验}
分别移除或替换各个模块，分析其对结果的影响，说明不同设计的作用。

\section{结论}
总结主要研究成果与应用意义，说明方法的局限，并提出未来研究方向。

\section*{致谢}
此处填写资助项目及致谢信息。

% 示例沿用 plainnat；投稿时按期刊要求选择参考文献样式。
\bibliographystyle{plainnat}
\bibliography{references}

\clearpage
\onecolumn
\appendix
\section{补充材料}
\label{app:details}
附录默认采用单栏排版，可放置详细推导、补充实验和实现细节。
如需保留双栏，删除本文件中的 \verb|\onecolumn|。

\subsection{复现信息}
记录数据版本、预处理步骤、超参数、随机种子以及软硬件环境。

\end{document}
